\section{An injective finite-group encoding}
\label{op:scalar-section}

We have a strict matrix one-sided inverse over $K[G_1]$. The
finite-factor reduction of Dykema and Juschenko
\cite[Lemma~2.1 and Theorem~2.2]{DykemaJuschenko2015} explains the
appropriate next step: put a full matrix algebra inside a finite group
algebra and fill the complementary summand with the identity. We prove
the actual matrix block needed here, including its injectivity. Its
shift and diagonal operators are the finite Heisenberg model; the
comparison with~\cite[Section~1.2]{GurevichHadani2009} is methodological,
since our cyclic order can be composite and our coefficient field is
finite.

\begin{proposition}\label{op:heisenberg}
Let $K$ be a field, let $M\ge2$ be invertible in $K$, and let
$\zeta\in K$ have order $M$. For an integer $q\ge0$ put
$Z_q=(\mathbb Z/M\mathbb Z)^q$ and $n=M^q$. There is a finite group
$\Lambda$ of order $Mn^2$ and a nonzero central idempotent
$f'\in K[\Lambda]$ for which explicitly specified elements
$\Omega_{uv}$, $u,v\in Z_q$, form a basis of $f'K[\Lambda]$ and obey
\[
 \Omega_{uv}\Omega_{st}=\delta_{vs}\Omega_{ut},\qquad
 \sum_u\Omega_{uu}=f'.
\]
Thus $\Mat_n(K)\cong f'K[\Lambda]$ by its standard matrix units.
\end{proposition}
\begin{proof}
Use the set $\Lambda=Z_q\times Z_q\times\mathbb Z/M\mathbb Z$ with
\begin{equation}\label{op:heisenberg-law}
 (x,y,k)(x',y',k')=(x+x',y+y',k+k'+y\cdot x').
\end{equation}
All coordinates and the dot product use residues modulo $M$.
Associativity follows from the bilinear identity
\[
 y\cdot x'+(y+y')\cdot x''=y'\cdot x''+y\cdot(x'+x'').
\]
The zero tuple is the identity, and the inverse is
$(-x,-y,-k+y\cdot x)$. No primality of $M$ is involved.
The element $z'=(0,0,1)$ is central of exact order $M$. Define
\[
 f'=\frac1M\sum_{k=0}^{M-1}\zeta^{-k}(z')^k.
\]
Reindexing gives $z'f'=\zeta f'$, and convolution gives $(f')^2=f'$.
The identity coefficient $M^{-1}$ is nonzero, so $f'\ne0$.

Write $X_x=(x,0,0)$ and $Y_y=(0,y,0)$, and put
\begin{equation}\label{op:heisenberg-D}
 D_v=\frac1n\sum_{y\in Z_q}\zeta^{-y\cdot v}f'Y_y.
\end{equation}
The denominator is invertible because $n=M^q$. Character orthogonality
here means the direct identity
\begin{equation}\label{op:composite-orthogonality}
 \sum_{y\in Z_q}\zeta^{y\cdot a}
 =\begin{cases}n,&a=0,\\0,&a\ne0.\end{cases}
\end{equation}
When $a\ne0$, choose a nonzero coordinate $a_j$. Translation of the
sum by the $j$th coordinate vector multiplies it by
$\zeta^{a_j}\ne1$. In the field $K$, this forces the sum to vanish.
The argument still works when $a_j$ is a zero divisor modulo the
composite integer $M$.

Convolution in the abelian group of $Y_y$ gives
$D_vD_t=\delta_{vt}D_t$; indeed the coefficient of $f'Y_b$ in the
product is
\[
 \frac{\zeta^{-b\cdot t}}{n^2}
       \sum_{y\in Z_q}\zeta^{y\cdot(t-v)}.
\]
Summing Equation~\eqref{op:heisenberg-D} over $v$ also gives
$\sum_vD_v=f'$. The group law gives
$Y_yX_a=(z')^{y\cdot a}X_aY_y$, and hence
\[
 D_vX_a=X_aD_{v-a}.
\]
Set
\begin{equation}\label{op:matrix-units}
 \Omega_{uv}=X_{u-v}D_v
 =X_{u-v}\frac1n\sum_{y\in Z_q}\zeta^{-y\cdot v}f'Y_y.
\end{equation}
Then
\[
 \Omega_{uv}\Omega_{st}
 =X_{u-v+s-t}D_{v-s+t}D_t
 =\delta_{vs}\Omega_{ut},
\]
and the diagonal sum is $f'$.

The multiplication identities alone do not yet prove that this block
has no kernel. On $K^{Z_q}$ define
\[
 \rho(x,y,k)\mathbf e_s=\zeta^{k+y\cdot s}\mathbf e_{s+x}.
\]
The scalar in a composite is
$\zeta^{k+k'+y\cdot x'+(y+y')\cdot s}$, so
Equation~\eqref{op:heisenberg-law} proves that $\rho$ is a representation.
It sends $f'$ to the identity and $D_v$ to the coordinate projection
at $v$. Therefore
\[
 \rho(\Omega_{uv})\mathbf e_s=\delta_{vs}\mathbf e_u,
\]
which proves independence of all $n^2$ elements $\Omega_{uv}$.
On the other hand, the $n^2$ elements $f'(x,y,0)$ span $f'K[\Lambda]$,
because $f'(x,y,k)=\zeta^k f'(x,y,0)$. Their group-basis supports have
distinct $(x,y)$ coordinates, so they too are independent. The sector
has dimension exactly $n^2$, and the matrix units form a basis.
When $q=0$, the empty tuple set has one element: $n=1$, $\Lambda=C_M$,
and the sole matrix unit is $f'$. All formulas remain valid.
\end{proof}

The elementary identity-completion step is worth recording separately.
Its unit need not be central in the target ring.

\begin{lemma}\label{op:complement}
Let $\Phi:\mathcal A\to\mathcal B$ be an injective ring homomorphism
between unital rings, not necessarily preserving their identities, and
put $g=\Phi(1_{\mathcal A})$. If $AB=1_{\mathcal A}$ and
$BA\ne1_{\mathcal A}$, then
\[
 a=1_{\mathcal B}-g+\Phi(A),\qquad
 b=1_{\mathcal B}-g+\Phi(B)
\]
satisfy $ab=1_{\mathcal B}$ and $ba\ne1_{\mathcal B}$.
\end{lemma}
\begin{proof}
The element $g$ is idempotent and is an identity on the image of $\Phi$
on both sides. Thus $1-g$ annihilates that image on both sides. Expanding
the two products gives
\[
 ab=1-g+\Phi(AB)=1,\qquad
 ba-1=\Phi(BA-1)\ne0.
\]
Only the last inequality uses injectivity; no global centrality of $g$
is required.
\end{proof}

\begin{corollary}\label{op:scalar}
Under the field hypotheses of Proposition~\ref{op:heisenberg}, suppose
$A,B\in\Mat_r(K[L])$ for a group $L$ and an integer $r\ge1$, with
$AB=I_r$ and $BA\ne I_r$. Finite formulas give a finite group $\Lambda$
and $a,b\in K[L\times\Lambda]$ with $ab=1$ and $ba\ne1$.
In particular, over any fixed finite field with at least three elements,
a failure of stable finiteness transfers to a scalar failure after a
finite direct factor, without changing the field.
\end{corollary}
\begin{proof}
Choose the least $q\ge0$ with $n=M^q\ge r$. This search is finite,
since $M^r\ge2^r\ge r$. Pad the matrices by identities:
\[
 \widetilde A=\diag(A,I_{n-r}),\qquad
 \widetilde B=\diag(B,I_{n-r}).
\]
The padding is empty if $n=r$. Then
$\widetilde A\widetilde B=I_n$ and
$\widetilde B\widetilde A-I_n=\diag(BA-I_r,0)\ne0$.
Use Proposition~\ref{op:heisenberg} and the group-basis identification
$K[L\times\Lambda]=K[L]\otimes_K K[\Lambda]$. For an $n$-square
matrix $T$ over $K[L]$, put
\[
 \Phi(T)=\sum_{u,v\in Z_q}T_{uv}\otimes\Omega_{uv}.
\]
The matrix-unit products prove multiplicativity, with
$\Phi(I_n)=1\otimes f'$. Independence of the $\Omega_{uv}$ gives
injectivity: extend their coordinate dual functionals to $K[\Lambda]$
and apply each in the second tensor factor to recover $T_{uv}$.
This works even when $K[L]$ has zero divisors.

Lemma~\ref{op:complement} now yields the promised finite sums
\begin{equation}\label{op:scalar-formulas}
 \begin{split}
 a&=1\otimes(1-f')+
       \sum_{u,v\in Z_q}\widetilde A_{uv}\otimes\Omega_{uv},\\
 b&=1\otimes(1-f')+
       \sum_{u,v\in Z_q}\widetilde B_{uv}\otimes\Omega_{uv}.
 \end{split}
\end{equation}
Their full scalar product is one, and their reverse defect is
$\Phi(\widetilde B\widetilde A-I_n)\ne0$.
For a finite field of order $Q\ge3$, take $M=Q-1$ and a generator
of its multiplicative group. Then $M\ge2$ is invertible in the field,
so the preceding construction applies.
\end{proof}

\begin{proof}[Proof of Theorem~\ref{op:main}]
Lemma~\ref{op:field} supplies the specified odd prime and field of
order $p^4$. Propositions~\ref{op:graph} and~\ref{op:modular} give the
graph fit and the strict regular-rank comparison.
Proposition~\ref{prop:first-gluing-injective} embeds the local groups
in $G_0$, and Proposition~\ref{prop:sector-ranks} places their ranks in
the nonzero ring $fK[G_1]$ with the required rectangular identities.
Corollary~\ref{op:edge-zero} supplies every ordered edge product.
Lemmas~\ref{op:compression}
and~\ref{op:kernel} produce $A_R,B_R\in\Mat_r(K[G_1])$, with their
nonzero certificate. Corollary~\ref{op:scalar}, using $L=G_1$ and
$M=p^4-1$, gives $G=G_1\times\Lambda$ and the sums
in Equation~\eqref{op:scalar-formulas}. Their product is one and their
reverse-product defect is nonzero by the injectivity of the matrix
encoding.

The group $G_0$ is generated by the finite groups $H_i$. The construction
of $G_1$ adds only $\langle z\rangle$ and the finitely many finite
groups $L_i$, and $\Lambda$ is finite. Thus $G$ is finitely generated.
The embedded copy of any $H_i$ supplies nontrivial torsion. All field,
group and coefficient choices were given by finite formulas and fixed
orders, as required.
\end{proof}

\paragraph{Finite prescriptions and scope.}
Every choice of character, basis, field element, or finite inverse
occurs in an explicitly finite set. Products of the resulting group
words define finite group-algebra sums even when different words happen
to coincide. No word comparisons are needed to define the elements,
and the kernel certificate and injective matrix block certify the
nonzero reverse defect. The countable permutation actions used to
certify group embeddings are not coefficient-search inputs.

Containing torsion does not say that every element has finite order.
The group must
be infinite, since finite-dimensional left multiplication proves direct
finiteness for every finite group algebra, in modular characteristic as
well as otherwise. The coefficient field throughout remains the
specified $K$ of order $p^4$.

More generally, the supports of any two finite sums in a group algebra
generate a finitely generated subgroup. Its group algebra embeds by the
group bases, preserving both a one-sided product and its nonzero reverse
defect. Thus a universal direct-finiteness theorem for finitely generated
groups would already imply the assertion for all groups.

\section{The cellular automaton on all configurations}
\label{op:cellular-section}

The final transfer is short and characteristic-independent. We give it
locally to fix the multiplication order and the meaning of a missing
configuration. Its group-ring ancestry appears in
\cite[Section~1]{ElekSzabo2004} and the linear theory
of~\cite{CSCCoornaert2006}.

For any field $k$, group $L$, and finite sum $c=\sum_uc_u u\in k[L]$,
define on all functions $x:L\to k$
\[
 (T_cx)(g)=\sum_uc_u x(gu).
\]
Only the coefficient sum is finite; $x$ may have infinite support.
Expanding the two finite sums gives
\begin{equation}\label{op:cellular-product}
 (T_cT_dx)(g)=\sum_{u,v}c_ud_vx(guv)=(T_{cd}x)(g).
\end{equation}
Thus $T_cT_d=T_{cd}$ and $T_1=\id$. If $\delta_h$ is the point
indicator at $h$, then $(T_c\delta_h)(1_L)=c_h$, which also proves
faithfulness.

\begin{proposition}\label{op:cellular}
Let $k$ be a finite field, $L$ a group, and $a,b\in k[L]$ with $ab=1$.
Then $T_b:k^L\to k^L$ is an injective cellular automaton, with
\[
 \im T_b=\{y\in k^L:T_{ba}y=y\}.
\]
For every $h$ with $(ba-1)_h\ne0$, the configuration $\delta_h$ is
absent from its image.
\end{proposition}
\begin{proof}
The support of $b$ is a finite memory set with rule
$(x_u)_u\mapsto\sum_u b_ux_u$. Equation~\eqref{op:cellular-product}
gives $T_aT_b=T_{ab}=\id$, proving injectivity. If $y=T_bx$, then
$T_{ba}y=T_bT_aT_bx=y$. Conversely, a fixed point of $T_{ba}$ has
$y=T_b(T_ay)$, hence lies in the image. Finally,
\[
 (T_{ba}\delta_h-\delta_h)(1_L)=(ba-1)_h\ne0,
\]
so $\delta_h$ is not a fixed point and cannot be in the image.
\end{proof}

Apply the proposition to the scalar sums in
Equation~\eqref{op:scalar-formulas}. Their nonzero defect has a
nonzero coefficient in the group basis, so the resulting automaton is
injective and nonsurjective. The missing point indicator has no
preimage even among infinitely supported configurations. If the finite
formula for $b$ displays a group element more than once, its terms
simply add in the same local rule; specifying the rule requires no
word-problem oracle.
Thus the constructed group also gives a negative answer to Gottschalk's
surjunctivity conjecture.
